\begin{table}[htbp]\centering
{\fontsize{10}{12}\selectfont\renewcommand{\arraystretch}{1.18}\setlength{\tabcolsep}{4pt}
\begin{tabularx}{\linewidth}{@{}C{24mm}L{112mm}X@{}}\toprule\rowcolor{tablehead}
\textbf{Notebook} & \textbf{Research stage} & \textbf{Code cells} \\ \midrule
00 & Questions, contracts and research design & 3 \\ \arrayrulecolor{tablerule}\hline
01 & Public API data and provenance & 4 \\ \arrayrulecolor{tablerule}\hline
02 & Order books, fees and liquidity & 4 \\ \arrayrulecolor{tablerule}\hline
03 & Payoff constraints and arbitrage & 4 \\ \arrayrulecolor{tablerule}\hline
04 & Probability modelling and model selection & 4 \\ \arrayrulecolor{tablerule}\hline
05 & Chronological validation and model-based value & 4 \\ \arrayrulecolor{tablerule}\hline
06 & Execution risk and research reflections & 4 \\ \arrayrulecolor{tablerule}\hline
07 & Synthesis and publication registers & 4 \\ \arrayrulecolor{tablerule}\hline
\arrayrulecolor{black}\bottomrule\end{tabularx}}
\caption[Executed notebook sequence]{Executed notebook sequence.}\label{tab:notebooks}
\vspace{2pt}{\centering\fontsize{9.3}{11.5}\selectfont Each notebook retains executed outputs. The HTML review groups them by milestone and provides individual downloads.\par}
\end{table}
